\cotefile{34}
\scanpage{133}
\mspage{V.5}
% archive sheet number: 132
% top diagram
\diagram{\begin{tikzcd}[ampersand replacement=\&,column sep=small,row sep=large]
\& X' \arrow[dl] \& X'_{\unsure{\eta|y}} \arrow[l] \arrow[dl] \& U' \arrow[l] \arrow[dl] \& U'_{\bar\eta} \arrow[l,no head] \arrow[dl] \\
X \arrow[d] \& X_\eta \arrow[l] \arrow[r] \arrow[d,"(\text{pas cart.})"] \& U \arrow[d] \& U_{\bar\eta} \arrow[l,no head] \arrow[d,no head] \\
S \& \eta \arrow[l] \& \eta \arrow[l,equal] \& \bar\eta \arrow[l,no head]
\end{tikzcd}}
% manuscript lines 1-3
1.4. Notons que l'on a
\[
X'=\varprojlim_\alpha X_\alpha,
\]
où les $X_\alpha$ sont étales sur $X$, et affines. On en conclut
\[
U'_{\deleted{\unsure{\bar\eta}}}=\varprojlim_\alpha X_{\alpha\deleted{\unsure{\bar\eta}}}\times_X U,\qquad
U'_{\bar\eta}=\varprojlim_\alpha(X_\alpha\times_X U)_{\bar\eta}
\]
% manuscript lines 4-8
d'où
\[
(\Psi^{(i)}_{\unsure{\nu|v}})_{\bar x}=H^i(U'_{\bar\eta},\Lambda)\simeq
\varinjlim_\alpha H^i((X_{\alpha\deleted{\unsure{\bar\eta}}}\times_X U)_{\bar\eta},\Lambda),
\tag*{(1.4.1)}
\]
où les \deleted{\unsure{$X_{\alpha\bar\eta}$ sont}} $(X_\alpha\times_X U)_{\bar\eta}$ sont des schémas \unsure{algébriques} quasi-affines sur $\bar\eta$, déduits par changement de \unsure{base} des schémas quasi-affines $X_\alpha\times_X U\subset X_\alpha$ sur $\eta$.

% manuscript lines 9-13
Lorsque $U=X_\eta$, cela devient

\deleted{2. Applications}
\[
\deleted{\unsure{\text{alors}}}\quad
(\Psi^{(i)}_\nu)_{\bar x}\unsure{\longrightarrow|=}H^i(\deleted{\unsure{X}}X'_{\bar\eta},\Lambda)\simeq
\varinjlim_\alpha H^i(X_{\alpha\bar\eta},\Lambda),
\tag*{(1.4.2)}
\]
où les $X_{\alpha\bar\eta}$ sont des schémas \inserted{alg.} \uline{affines} sur $\bar\eta$ [déduits des schémas affines $X_{\alpha\eta}$ sur $\eta$]. \deleted{\unsure{les}}

% manuscript lines 14-16
NB. Plus généralement, dans 1.1.1. les $(X_\alpha\times_X U)_{\bar\eta}$ sont affines si on suppose que $U\mathbin{\unsure{-|\to}}X$ est affine, i.e. $U\to X_\eta$ est affine.
\scanpage{134}
\mspage{V.6}
% archive sheet number: 133
% manuscript lines 1-3
2. \uline{Application du th. de Lefschetz affine.}

2.1. Soit \deleted{\unsure{$a$}} $x\in X_0$, et
\[
n_x=\deleted{\dim X_{\unsure{\eta|y}}\unsure{\text{ local}}}
\sup_{\substack{X_i\text{ comp. irréduct.}\\\text{de }X\text{ passant par }x}}\dim X_{i\eta}
\quad\deleted{\unsure{i\in I}}
\tag*{(2.1.0)}
\]
\inserted{Supposons $U\to X_\eta$ affine.}
% manuscript lines 4-7
Alors dans (1.4.2) les $X_{\alpha\bar\eta}$ sont des affines et \unsure{finis|quasi-}\deleted{\unsure{compacts $U_{\bar\eta}$}} étales sur $X_\eta$, \inserted{et si on veut au dessus de $\bigcup_i X_{i\eta}$;} donc affines de dim $\leq n$, donc leur $H^i$ est nul pour $i>n$ \inserted{(SGA \deleted{\unsure{4}}4 \uline{XIV} ---)}. On trouve donc
\[
(\Psi^i_{\unsure{\nu|v|*}})_{\bar x}=0\qquad\text{pour }i>n_x.
\tag*{(2.1.\unsure{1})}
\]
% manuscript lines 8-11
On va améliorer ce résultat en prouvant que
\[
(\Psi^i_\nu)_{\bar x}=0\qquad\text{pour }i>n_x-\dim\overline{\{x\}}.
\tag*{(2.1.2)}
\]
ce qui implique aussi
\[
(\Psi^i_\nu)_{\bar x}=0\qquad\text{si }i>\dim\underline{\unsure{\mathcal O|\mathscr O}}_{X_0,x}=\dim_x(X_0)-\dim\overline{\{x\}}
\tag*{(2.1.3)}
\]
% manuscript lines 12-18
2.2. Pour ceci, notons qu'on peut \deleted{\unsure{supposer (sans changer}} remplacer $X$ par \deleted{$X_{\unsure{\bar\eta}}$} l'adhérence schématique $\overline{(X_\eta)_{\unsure{\mathrm{réd}|\mathrm{red}}}}$ [sans changer les \deleted{\unsure{valeurs des}} $\Psi^{(i)}_\nu$ \deleted{\unsure{\}}} ni $n_x$, et en diminuant $\dim_x X_0$], alors \deleted{\unsure{les deux} $n_x-\dim_x X_0$, \unsure{et les deux}} \deleted{\unsure{de}} $n_x=\dim_x X_0$ \deleted{$\dim X_0\leq\unsure{n}$} et les deux formules (2.1.1) et (2.1.2) sont identiques.

\deleted{\unsure{quitte à remplacer $X$ par un voisinage ouvert affine de $x$ dans $X$,}}

% manuscript lines 19-25: entire block cancelled by long diagonal strokes
\deleted{Alors, il existe un morphisme fini}
\[
\deleted{X\xrightarrow{h}Y=S[T_1,\ldots,T_n],}
\]
\deleted{et si on pose}
\[
\deleted{F=h_*(\Lambda_X),\quad y=h(x),}
\]
\deleted{on trouve que les faisceaux $\Psi^{(i)}_{\unsure{f|g}}(F)$ relatifs à $F$ (qui se définissent de façon évidente, comme}

% manuscript lines 26-30
On peut supposer aussi (quitte à \unsure{restreindre|rétrécir} $X$) que \deleted{$n_x=\dim X_\eta$; \unsure{Posons, $n_x$} \unsure{et $n$}. Soit}
\[
\deleted{\unsure{n=\dim}}\qquad n=\dim X_\eta=n_x;
\]
Soit d'autre part
\[
r=\text{degtr. }k(x)/k.
\]
\scanpage{135}
\mspage{V.7}
% archive sheet number: 134
% manuscript lines 1-6
\deleted{\unsure{et soit}}

On peut supposer $X$ affine, et on peut trouver une suite d'\unsure{élts|éléments}
\[
t_1,\ldots,t_r\in\Gamma(X,\unsure{\mathcal O|O}_X)
\]
telle que leurs images dans $k(x)$ \unsure{font|forment} une base de transc. sur $k$. Les $t_i$ définissent un morphisme
\[
h:X\mathbin{\unsure{-|\to}}Y=S[T_1,\ldots,T_r],
\]
% manuscript lines 7-11
et si $y=h(x)$, on voit que $k(x)/k(y)$ est fini, donc que $x$ est fermé dans $h^{-1}(y)$. On a donc
\[
\dim_x h^{-1}(y)=\dim\underline{\mathcal O}_{X_0,x}=n-r,
\]
et quitte à restreindre $X$, on peut donc supposer que les fibres de $X$ sur $Y$ sont de dim $\leq n-r$.

% manuscript lines 12-15
Soit alors $S_1$ le localisé strict de $Y$ en le \deleted{le pt fermé} pt géométrique de $Y$ défini par le pt géom $\bar x$ de $X$, $X_1=X\times_Y S_1$, $U_1=U\times_Y S_1$. Alors $X_1$ est affine

% lower-left diagram: entire diagram cancelled by diagonal strokes
\diagram{\begin{tikzcd}[ampersand replacement=\&,column sep=small,row sep=small]
\&\deleted{X'}\arrow[d]\&\deleted{U'}\arrow[l,no head]\arrow[d]\&\deleted{U'_{\unsure{\bar\eta|\bar y}}}\arrow[d,no head]\&\deleted{\unsure{U'\otimes_{S_1}\bar L}}\arrow[l,no head]\arrow[d,no head]\\
\deleted{X}\arrow[d]\&\deleted{X_1}\arrow[l]\arrow[d]\&\deleted{\unsure{U_1}}\arrow[d,"\deleted{\text{\unsure{pas cart.}}}"]\\
\deleted{S[T_1,\ldots,T_r]=Y}\arrow[d]\&\deleted{S_1}\arrow[l]\arrow[d,"\deleted{\unsure{\varphi}}"]\&\deleted{S_1}\arrow[l,equal]\arrow[d,no head]\&\deleted{S_{1\bar\eta}}\arrow[l,no head]\arrow[d,no head]\&\deleted{\bar L}\arrow[l,no head]\\
\deleted{S}\&\deleted{S}\&\deleted{S}\arrow[l,no head]\&\deleted{\bar\eta}\arrow[l,no head]
\end{tikzcd}}
% lower-right revised diagram
\diagram{\begin{tikzcd}[ampersand replacement=\&,column sep=small,row sep=small]
\&X'\arrow[d]\&U'\arrow[l]\arrow[d,no head]\&U'_{\bar\eta}\arrow[l]\arrow[dd,"\varphi"]\&U'\times_{S_1}\bar\lambda\arrow[l,no head]\arrow[dd]\\
X\arrow[d]\&X_1\arrow[l]\arrow[d]\&U_1\arrow[l]\arrow[d]\\
S[T_1,\ldots,T_r]=Y\arrow[d]\&S_1\arrow[l]\&S_1\arrow[l,equal]\arrow[d]\&S_{1\bar\eta}\deleted{\unsure{=\lambda}}\arrow[l]\arrow[d,no head]\&\deleted{\unsure{\bar L}}\bar\lambda\arrow[l,no head]\\
S\&\deleted{\unsure{\eta}}\arrow[u]\&S\&\deleted{\unsure{\eta}}\bar\eta\arrow[l]\&\deleted{\unsure{\longrightarrow|-}}
\end{tikzcd}}
% manuscript lines 16-20
\unsure{t.f.|l.f.} sur le trait \unsure{st.} hensélien $S_1$, de dim relative $n-r$, $U_1$ en est un ouvert affine, $X'$ est le localisé \inserted{strict} \uline{de} $X_1$ en le pt géom $\bar x$, $U'$ est l'image inverse de $U_1$ dans $X'$. Soit $L$ le corps des fractions de $S_1$ qui est un
\scanpage{136}
\mspage{V.8}
% archive sheet number: 135
% manuscript lines 1-4
extension \inserted{algébrique} séparable de $K(T_1,\ldots,T_r)$, donc une ext. séparable de $K=k(\unsure{\eta|y})$, (et $\lambda=\operatorname{Spec}L\simeq(S_1)_\eta$ \deleted{\unsure{est un}}). D'après 2.1.1. appliqué à la situation $X_1\to S_1$, et $U_1\subset X_1$, $X'\mathbin{\unsure{-|\to}}X_1$, on voit que
\[
H^i(U'\times_{S_1}\bar\lambda,\Lambda)=0\quad\text{si }i>n-r,\qquad
\text{i.e. }R^i\varphi_*(\Lambda_{U'_{\bar\eta}})=0\quad\text{si }i>n-r
\tag*{(2.\unsure{2|1}.1)}
\]
% manuscript lines 5-8
D'autre part, \deleted{$\bar\lambda$ \unsure{est un}} on sait que \deleted{\unsure{$S_1$}} ($S_1\mathbin{\unsure{-|\to}}S$ étant déduit par localisation stricte du morphisme lisse $Y\to S$) que $S_{1\eta}=\lambda$ est acyclique pour les nos premiers $\ne p$ (SGA 4 XV). Donc la suite spectrale
% slanted
\margin{réf. inutile, on voit directement que $\operatorname{Gal}(\unsure{\bar\lambda|\bar\eta}/S_{1\eta})$ est un pro-$p$-groupe\unsure{---}.}
\[
H^i(U'_{\bar\eta},\Lambda)\Longleftarrow\mathbb E_{\unsure{2|,}}^{p,q}=H^p(\lambda,R^q\varphi_*(\Lambda_{U'_{\bar\eta}}))
\]
% manuscript lines 9-13
dégénère et donne
\[
H^i(U'_{\bar\eta},\Lambda)\simeq H^0(\lambda,R^i\varphi_*(\Lambda_{U'_{\bar\eta}}))\hookrightarrow H^i(U'\times_{S_1}\bar\lambda,\Lambda)
\tag*{(2.2.2)}
\]
ce qui implique par (2.2.1)
\[
H^i(U'_{\bar\eta},\Lambda)=0\qquad\text{pour }i>n-r,
\]
et achève la démonstration. On a prouvé :

% manuscript lines 14-17
\uline{Théorème} 2.3. Avec les notations de 2.1, on a \inserted{pour $x\in X_0$,}
\[
(\Psi^q_{\unsure{\nu|v}})_{\bar x}=0\qquad\text{pour }q>n_x-\dim\overline{\{x\}}
\]
(et à fortiori pour $i>\dim\underline{\unsure{\mathcal O|O}}_{X_0,x}=\dim_x X_{\unsure{x|0}}-\dim\overline{\{x\}}$).
De façon équivalente on a
\[
(\underline\Phi^q_\nu)_{\bar x}=0\quad\text{si }q>n_x+1-\dim\overline{\{x\}},\quad q\geq1.
\]
% manuscript lines 18-23: cancelled block
\deleted{Appliquant le théorème de dim. cohomologique SGA 4 \uline{IX}, on trouve}
\deleted{\inserted{si $n=\dim X_\eta$, $q\geq1$, $X=\overline{X_\eta}$ donc $\underline\Phi^0_\nu=0$ :}}
\[
\deleted{H^p(X_0,(\underline\Phi^q_\nu)|X_0)=0\quad\text{si }p\geq2(n+1-q),}
\]
\deleted{i.e. $p+q\geq2n+2-q$}

\deleted{\uline{Corollaire} 2.\unsure{3}. On a, si $X$ propre sur $S$}

% manuscript lines 24-26
\uline{Corollaire} 2.\unsure{4|3}. On a : $\underline\Psi^q_\nu=0$ si $q>\dim X_\eta$, $\underline\Phi^q_\nu=0$ si $q>\dim X_\eta+1$.
\scanpage{137}
\mspage{V.9}
% archive sheet number: 136
% manuscript lines 1-4
\deleted{\unsure{2.4.}} \deleted{Supposons} \inserted{Considérons} maintenant \deleted{que} l'ens \inserted{(fermé) $Z$} des points de $X_0$ où $f:X\mathbin{\unsure{-|\to}}S$ est non lisse, \deleted{\unsure{soit de}} \inserted{$d$ sa} dim. \deleted{\unsure{$d$}}
On a \deleted{alors} \inserted{$\operatorname{supp}\underline\Phi^q_{\unsure{\nu|v}}\subset Z$ (\uline{III} 1.), donc,} par SGA 4 \deleted{\unsure{6}} \uline{X}
\[
H^p(X_0,\underline\Phi^q_\nu)\ne0\quad\deleted{\text{\unsure{signifie}}}\inserted{\text{implique}}\quad p\leq2d.
\tag*{(*)}
\]
% manuscript lines 5-9
D'ailleurs, on a en vertu de 2.\unsure{3}. (où $n=\dim X_{\unsure{\bar\eta|\bar y}}$)
\[
\delta(\underline\Phi^q_\nu)\leq n+1-q\qquad\text{si }q\geq1
\]
donc on trouve une \unsure{valeur|résultat} \inserted{$\delta(\Phi^q_\nu)$} $\leq d$ si $n+1-q<d$ i.e. $q>n+1-d$; \deleted{donc} pour de tels $q$, on a par SGA 4 \uline{X}
\[
\begin{aligned}
&\deleted{2\delta({\unsure{\underline\Phi|\Phi}}^q_\nu)\unsure{.}}\\
&H^p(X_0,\underline\Phi^q_\nu)\ne0\quad\deleted{\text{\unsure{signifie}}}\inserted{\text{implique}}\quad p\leq2(n+1-q)\\
&\hspace{8em}\deleted{\text{i.e. implique}}\qquad p+q\leq2n+2-q,
\end{aligned}
\tag*{(**)}
\]
% manuscript lines 10-14
et on a \inserted{pour ces $q$,} $2n+2-q<2n+2-(n+1-d)=n+1+d$. \unsure{Donc}
D'autre part, si $q\leq n+1-d$ et \inserted{$H^p(X_0,\underline\Phi^q_\nu)\ne0$ i.e.} $p\leq2d$, on a $p+q\leq n+1+d$. Donc en tous cas (si $q\geq1$)
\[
H^p(X_0,\underline\Phi^q_\nu)\ne0\quad\Longrightarrow\quad p+q\leq n+d+1.
\]
% manuscript lines 15-19
\deleted{Supposons} Cette implication est aussi vraie pour $q=0$, si on suppose que $d\leq n+1$ (ce qui est le cas \inserted{(et $d\leq n=\dim X_0$), p. ex. si $X=\overline{X_\eta}$,} \deleted{\unsure{en particulier}} \unsure{i.e.} $X$ plat sur $S$]. On trouve \inserted{alors (\unsure{cf. 1})}
\[
\Phi^i_\nu=0\qquad\text{pour }i>n+d+1.
\]
\inserted{Donc :}

% manuscript lines 20-27
\uline{Corollaire} 2.5. \deleted{Supposons que} Soit $d$ la dimension de l'ens des pts de $X_0$ en lesquels $f:X\mathbin{\unsure{-|\to}}S$ est non lisse, \deleted{et} \inserted{$n=\dim X_{\bar\eta}$, et} supposons que $d\leq n+1$ i.e. que $\dim X_0\leq n+1$. [ce qui est le cas si $X=\overline{X_\eta}$, impliquant $\dim X_0=n$]. Alors le \unsure{grp|grps} $\Phi^i_\nu$ des cycles évanescents \inserted{(\uline{III} 1.)}, à coefficients dans $\Lambda_\nu=\mathbb Z/\ell^\nu\mathbb Z$, est nul pour \deleted{\unsure{$q$}} $i>n+d+1$, en d'autres termes l'hom
% slanted
\margin{! (rel. à $U=X_\eta$)}
\scanpage{138}
\mspage{V.10}
% archive sheet number: 137
% manuscript lines 1-4
\[
H^i(X_0,\Lambda_{\unsure{\nu|v}})\longrightarrow H^i(X_{\unsure{\bar\eta|\bar y}},\Lambda_\nu)
\]
est un \unsure{isomorphisme} si $i>n+d+1$, un épimorphisme si $i=n+d+1$. L'action du groupe d'inertie \deleted{sur} $I$ sur $H^i(X_{\bar\eta},\Lambda_\nu)$ est triviale si $i\geq n+d+1$.

% manuscript lines 5-10
Procédant par dualité on trouve :

\uline{Corollaire} 2.6. Avec les hypothèses précédentes, l'action de $I$ sur $H^i(X_{\bar\eta},\Lambda_\nu)$ est triviale si $i\notin[n-d,n+d]$.
En particulier, si $d=0$ i.e. si l'ens $Z$ des pts de $X_0$ où $f$ est non lisse est fini, alors l'action de $I$ sur $H^i(X_{\bar\eta},\Lambda_\nu)$ est triviale si $i\ne n$.

% manuscript lines 11-15
\uline{Remarque} 2.7. Lorsque $f:X\mathbin{\unsure{-|\to}}S$ est projectif, on peut trouver une démonstration plus simple de \deleted{\unsure{trivialité}} 2.5 (et \inserted{de} 2.6) \deleted{\unsure{de l'action triviale du groupe d'inertie}}. En effet, \deleted{\unsure{plongeons}} utilisons un plongement $X\subset\mathbb P^r_S$, et soit $Y$ un sous-schéma de $X$, obtenu en coupant $X$ avec $d+1$ \deleted{\unsure{sections}} hyperplans de $\mathbb P^r_S$.
% slanted
\margin{intervertir avec 2.8}
% manuscript lines 16-22
Si ces derniers sont bien choisis, on peut supposer $Y_0$ lisse de codim $d+1$ dans $X_0$, et alors $Y$ est lisse sur $S$, de codim $d$ dans $X$. De plus, le th. de Lefschetz affine SGA 4 XIV implique aisément, \deleted{\unsure{par les}} de façon connue, que l'hom de Gysin
\[
H^i(Y_{\bar\eta},\Lambda_\nu)\mathbin{-}H^{\deleted{\unsure{i}}\inserted{i+2(d+1)}}(X_{\bar\eta},\Lambda_\nu(\deleted{\unsure{+|2|*}}d+1))
\]
\scanpage{139}
\mspage{V.11}
% archive sheet number: 138
% manuscript lines 1-3
est surjectif si $i\geq\dim Y_{\unsure{\bar\eta|\bar y}}=n-d-1$, \deleted{\unsure{bi|ou}} donc si $i+2(d+1)\geq n+d+1$, \inserted{et bijectif si $i>n-d-1$ i.e. si $i+2(d+1)>n+d+1$,}
% manuscript lines 4-7: cancelled argument
\deleted{Donc le fait que $I$ opère trivialement sur $H^i(X_{\bar\eta},\Lambda_{\unsure{\nu|v}})$ est acquis puisque le fait qu'il opère trivialement sur $H^i(Y_{\bar\eta},\Lambda_\nu)$, qui est \unsure{une} conséquence bien connue du fait que $Y$ est propre et lisse sur $S$ (SGA 4 XVI ---).}
\inserted{Le diagramme commutatif}
\inserted{\diagram{\begin{tikzcd}[ampersand replacement=\&]
H^i(\deleted{\unsure{X}}Y_{\bar s},\Lambda)\arrow[r]\arrow[d,"\varphi_Y"']\&H^{i+2(d+1)}(X_0)\arrow[d,"\varphi_X"]\\
H^i(Y_{\bar\eta},\Lambda)\arrow[r]\&H^{i+2(d+1)}(\deleted{\unsure{Y}}X_{\bar\eta})
\end{tikzcd}}}
% slanted
\margin{donne alors le résultat voulu ($\varphi_Y$ est bijectif par SGA 4 \uline{XVI} ---)}

% manuscript lines 8-12
\deleted{\unsure{2.}}\uline{Remarque} 2.8. L'assertion 2.3 suggère la \deleted{validité de} la question suivante. Soit $X$ \inserted{$=\operatorname{Spec}A$} un schéma strictement local \inserted{de dim. \deleted{\unsure{un}} $n$}, $t\in\unsure{\mathfrak m|\mathfrak r}_A$, \deleted{un élément de et} $U=X-V(t)=\operatorname{Spec}A_t$, et soit pour les entiers $m$ premiers à la car. résiduelle
\[
\unsure{X_m|X_{(m)}}=\operatorname{Spec}A[T]/(T^m-t);\qquad U(m)=X_m\times_X U,
\]
enfin soit
\[
\deleted{\unsure{X}}U(\infty)=\varprojlim U(m).
\]
% manuscript lines 13-18
Est-il vrai que la dim. cohomologique de $U(\infty)$, pour des coefficients premiers à \deleted{la} $p$, est $\leq n-1$ ? [La conjecture standard SGA 4 XIV \unsure{et|e} \deleted{donnerait par} donnerait \inserted{pour $\ell\ne p$ :} $\operatorname{cd}_\ell(U(m))\leq\unsure{n}$ pour $\forall\,m$, donc $\operatorname{cd}_\ell(U(\infty))\leq n$, et non $n-1$ !
\scanpage{140}
\mspage{V.12}
% archive sheet number: 139
% manuscript heading
3. \uline{Application du th. de Lefschetz \deleted{\illegible{1}} \inserted{(de la section hyperplane)} \unsure{global}\deleted{\unsure{e}}.}

% manuscript lines 1-5
\deleted{\unsure{3.1|3.4}} 3.1. Supposons que la \uline{profondeur étale rectifiée} \inserted{pour $\ell$} de $X_s$ et $X_{\unsure{\eta|y}}$ soit $\geq m$, i.e. que pour tout $x\in X_s$, on ait
\[
\overbrace{\operatorname{prof\,et}_x X_s+\dim\overline{\{x\}}}^{\inserted{\operatorname{prof\,et}_x^{*}(X_s)}}\geq m
\tag*{(3.1.1)}
\]
et de même pour $\forall\ x\in X_\eta$,
\[
\underbrace{\operatorname{prof\,et}_x(X_\eta)+\dim\overline{\{x\}}}_{\inserted{\operatorname{prof\,et}_x^{*}(X_\eta)}}\geq m,
\tag*{(3.1.2)}
\]
% manuscript lines 6-9
où ici $\overline{\{x\}}$ désigne l'adhérence dans $X_\eta$. (On note que l'hypothèse faite sur $X_\eta$ ne dépend en fait que de la fibre \uline{géométrique} $X_{\bar\eta}$).
% slanted
\margin{! Cf SGA 1 \uline{XIV} --- pour la notion de \uline{profondeur étale}.}
\deleted{Alors, \unsure{en}} Soit $d$ comme dans 2.5, et supposons \inserted{d'abord} $X$ projectif et \uline{plat sur} $S$.
% manuscript lines 10-15
Choisissons $Y$ comme dans 2.7, alors le th. de Lefschetz \deleted{\unsure{p}} \inserted{global} sous la forme SGA 1 \uline{XIV} 4.6 \uline{entraîne} \uline{[sous réserve de disposer de la résolution des singularités sur \deleted{\unsure{tous}} les \unsure{clôtures|clotures} alg. de $k$ et $K$.]} :
\[
H^i(X_s,\Lambda)\mathbin{-}H^i(Y_s,\Lambda)
\quad(\text{resp. }H^i(X_{\bar\eta},\Lambda)\mathbin{-}H^i(Y_{\bar\eta},\Lambda))
\]
est bijectif pour $i<m-d-1$, injectif pour $i=m-d-1$. Le diagramme commutatif
\diagram{\begin{tikzcd}[ampersand replacement=\&]
H^i(X_s,\Lambda)\arrow[r,"\varphi_X^i"]\arrow[d]\&H^i(X_{\bar\eta},\Lambda)\arrow[d]\\
H^i(Y_s,\Lambda)\arrow[r,"\varphi_Y^i","\sim"']\&H^i(Y_{\bar\eta},\Lambda)
\end{tikzcd}}
% manuscript lines 16-19
(où $\varphi_Y$ est bijectif en vertu de SGA 4 XVI ---) montre que \deleted{$\varphi_X^i$ est bijectif}
\[
\deleted{\unsure{\text{(3.1)}}}\qquad\varphi_X^i:H^i(X_s,\Lambda)\mathbin{-}H^i(X_{\bar\eta},\Lambda)
\]
\uline{est bijectif pour $i<m-d-1$, injectif pour} $i=m-d-1$,
\scanpage{141}
\mspage{V.13}
% archive sheet number: 140
% manuscript lines 1-2
en d'autres termes que l'on a
\[
\Phi^i_{\unsure{\nu|v}}=0\qquad\text{pour }i\leq m-d-1.
\tag*{(3.1.3)}
\]
% Everything below is crossed by long oblique cancellation strokes.
% manuscript lines 3-8
\deleted{Remarque 3.2. Supposons par exemple que $X_s$ \deleted{et $X_\eta$} soit \unsure{localement} \deleted{une} intersection complète, \deleted{\unsure{pour que $X$ soit}} \inserted{ce qui vient au même que $X$ soit loc. une intersection complète,} \deleted{\unsure{implique}} ce qui implique la même chose pour $X_\eta$. Alors (\deleted{\unsure{cf.}} SGA 1 \uline{XIV} 5.6) on peut prendre $m\geq n$, et (3.1.3) prend la forme}
% slanted; the margin note is itself crossed by a cancellation stroke
\margin{\deleted{Mais dans ce cas on peut améliorer la conclusion, en remplaçant $i\leq\unsure{n|m}-d-1$ par $i\leq n-d$.}}
\[
\deleted{\Phi^i_\nu=0\qquad\text{pour }i\leq n-d-1,}
\tag*{\deleted{(3.2.1)}}
\]
% manuscript lines 9-14
\deleted{ce qui, compte tenu de 2.7, prend la forme}
\[
\deleted{\Phi^i_\nu=0\qquad\text{si }i\notin[n-d,n+d+1]}
\tag*{\deleted{(\unsure{3.1.5|3.2.5})}}
\]
\deleted{Si $d=0$, \deleted{\unsure{par}} (cas d'un \deleted{\unsure{singularités}} lieu de non lissité \uline{fini} sur $X_0$) cela signifie}
\[
\deleted{\Phi^i_\nu=0\qquad\text{si }i\ne n,n+1.}
\tag*{\deleted{(3.1.6)}}
\]
% manuscript lines 15-24
\deleted{Du point de vue action du groupe d'inertie $I$, (3.1.3) \deleted{\unsure{signifie}} implique que cette action est triviale \deleted{sur} $H^i(X_{\unsure{\bar\eta|\bar y}},\Lambda_\nu)$ si $i<m-d-1$, donc compte tenu de 2.7) si $i\notin[m-d-1,n+d]$, donc si $m=n$ (p. ex. $X_s$ loc. une intersection complète) si $i\notin[n-d-1,n+d]$, et si de plus $d=0$ si \deleted{\unsure{$i\notin[n-1,n]$}} $i\ne n-1,n$. On comparera avec la conclusion de 2.6, plus fort (car il suffit de supposer $i\ne n$) mais avec une hypothèse \inserted{de régularité locale (\unsure{de nature assez exclusive})} différente [$X_\eta$ lisse au lieu de $X_s$ d'intersection complète].}
\scanpage{142}
\mspage{V.14}
% archive sheet number: 141
% manuscript lines 1-8
\uline{Remarque} 3.2. a) Lorsque $X_s$ est \inserted{localement} une intersection complète, \deleted{alors} et purement de dimension $n$, alors \deleted{\unsure{l'hypothèse}} \deleted{les hypothèses} il en est de même de $X_{\unsure{\eta|y}}$, et les hypothèses (3.1.1) et (3.1.2) sont vérifiées, avec $m=n$. Nous verrons cependant plus bas (au \deleted{\unsure{\S}}\S\ 5) que l'on peut alors (sous réserve de conjectures rappelées \deleted{dans (\unsure{2.2})} tout à l'heure 3.3)) améliorer le résultat \inserted{(3.1.3)} en écrivant $\underline\Phi^i_{\unsure{\nu|v}}=0$ pour $i\leq n-d$.

% manuscript lines 9-13
b) Par contre, \deleted{\unsure{bien}} même si $X_{\bar\eta}$ et $X_s$ sont \deleted{\unsure{irréductibles}} \inserted{purement} de dimension $n$, \deleted{et si $n=1$}, et si les hypothèses (3.1.1) et (3.1.2) sont satisfaites \deleted{pour} avec $m=n$, on ne peut en général améliorer la conclusion (3.1.3) d'une unité comme dans a).

% manuscript lines 14-23
\deleted{Pour} Pour le voir, \unsure{partons|prenons} p. ex. de \deleted{\unsure{un} $X'$ lisse sur $S$ à fibres géom. \deleted{\unsure{int.}} connexes (p. ex.} \deleted{$X'=$} $X'=\mathbb P^1_S\amalg\mathbb P^1_S$ et \deleted{\unsure{par}} un $X$ déduit de $X'$ en identifiant deux pts \inserted{distincts} de $X'_s$ rationnels sur $k$, \inserted{appartenant à deux composantes distinctes de $X'_s$}. Alors \deleted{les hypothèses} $X$ satisfait les hypothèses avec $m=n=1$, $d=0$, cependant \deleted{la conclusion} on n'a pas $\Phi^1=0$, \deleted{car \unsure{contrairement}} \inserted{i.e.} l'hom. canonique $H^0(X_0)\mathbin{-}H^{\unsure{0|1}}(X_{\bar\eta})$ n'est pas bijectif, puisque $X_0$ est connexe mais $X_{\bar\eta}$ ne l'est pas.
\scanpage{143}
\mspage{V.15}
% archive sheet number: 142
% manuscript lines 1-7
3.\unsure{3}. Notre but dans la suite de ce \S\ est de donner des résultats locaux sur la nullité \unsure{des|de} $(\underline\Phi^i_{\unsure{\nu|v}})_{\bar x}$ analogues aux résultats globaux précédents, qui impliquent ces derniers sans hypothèse \unsure{projectivité} pour $f:X\to S$. L'outil utilisé sera le th. \uline{de Lefschetz local} (SGA 2 \uline{XIV} 4.5), qui remplace le th. de Lefschetz global.

% manuscript lines 8-17
Il n'est prouvé que sous réserve de disposer des th. de pureté, \deleted{et} des th. de résolution des singularités et du th. de \uline{Lefschetz affine local}, \deleted{dans diverses} Compte tenu des résultats de SGA 4 \uline{XIV}\deleted{\unsure{4}}, \uline{XVI}, \uline{XIX}, il en est ainsi \deleted{\unsure{avec})} \deleted{\unsure{si}} \unsure{$S$} d'égale caractéristique \inserted{$p$}, lorsque on dispose de la résolution des \uline{singularités} pour les schémas excellents en cette caractéristique --- \deleted{\unsure{en part.} \deleted{\unsure{si $S$ de}} caractéristique nulle} donc pour le moment, seulement \uline{si} $S$ \uline{est de caractéristique nulle}.

% manuscript lines 18-23
Dans la suite, nous supposons que nous disposons des ingrédients que nous venons d'expliciter; cela nous oblige pour le moment à nous borner au cas de la car. nulle, ou au cas d'égale car. avec des restrictions dimensionnelles draconiennes ($\dim X_{\unsure{\eta|y}}\leq2$).
\scanpage{144}
\mspage{V.16}
% archive sheet number: 143

\deleted{3.4.}\inserted{4.} \deleted{Les cinq} \uline{Situation locale des cinq schémas.}

4.1 Soient \deleted{$\unsure{\bar x\longmapsto x}$} $x\in X_0$, $\bar x$ un pt géométrique sur $x$, $X'$ le localisé strict de $X$ en $\bar x$, posons
\[
\tag*{(4.1.1)}
V=X'-\{\bar x\},\quad V_s=V\times_S\{s\},\quad
V_{\unsure{\eta|y}}=V\times_S\eta=V-V_s,\quad
V_{\bar\eta}=V\times_S\bar\eta=V_\eta\times_\eta\bar\eta.
\]
En plus de ces quatre schémas remarquables, nous en introduisons un cinquième. \deleted{à} \inserted{Pour \unsure{cela}, considérons} \deleted{\unsure{le schéma}} le normalisé $\bar S$ de $S$ dans $\bar K$ :
\[
\tag*{(4.1.2)}
\bar S\longrightarrow S,\qquad \bar S=\varprojlim S_\alpha
\]
où les $S_\alpha$ sont les normalisés de $S$ dans les extensions finies \deleted{\unsure{du point}} $K_\alpha$ de $\bar K/K$. \deleted{+} On posera
\[
\tag*{(4.1.3)}
\bar V=V\times_S\bar S=\varprojlim V_\alpha,\qquad V_\alpha=V\times_S S_\alpha.
\]
\deleted{\unsure{Soit alors} $\bar V_s=\unsure{V_s'}$.}
Evidemment on a un isomorphisme
\[
\deleted{V_{\bar\eta}\overset{\unsure{\mathrm{dfn}|\mathrm{df}}}{=}\bar V\times_S\eta
\simeq\bar V_\eta\overset{\mathrm{dfn}}{=}\bar V\times_{\bar S}\bar\eta\simeq V_{\bar\eta}}
\]
\margin{\diagram{\begin{tikzcd}[ampersand replacement=\&,column sep=small,row sep=small]
\& V\arrow[ddd] \& \& \bar V\arrow[ll]\arrow[ddd] \& \\
\& \& V_\eta\arrow[ul]\arrow[rr]\arrow[ddd] \& \& V_{\bar\eta}\arrow[ul]\arrow[ddd] \\
s\arrow[dr] \& \& \& \bar s\arrow[d] \& \\
\& S \& \& \bar S\arrow[ll] \& \\
\& \& \eta\arrow[ul] \& \& \bar\eta\arrow[ll]\arrow[ul]
\end{tikzcd}}}
\margin{\diagram{\begin{tikzcd}[ampersand replacement=\&,column sep=small,row sep=small]
\deleted{V}\arrow[r]\arrow[d] \& \deleted{\bar V}\arrow[d] \\
\deleted{S}\arrow[r] \& \deleted{\bar S}
\end{tikzcd}}}
\deleted{$V_{\bar\eta}$}
\[
\tag*{(4.1.4)}
V_{\bar\eta}\lhook\joinrel\longrightarrow\bar V,
\]
\unsure{identifiant|identifie} $V_{\bar\eta}$ à la fibre générique de $\bar V$ sur $S$ [resp. sur $\bar S$], et le complémentaire \deleted{$\unsure{V_\eta}$} $\bar V-V_{\bar\eta}$ \deleted{\unsure{est aussi} $\unsure{V^{\prime\prime}}$} \inserted{s'identifie à} \deleted{$V_s\unsure{\times_S S}$}
\[
\tag*{(4.1.5).}
V_{\bar s}=V_s\otimes_{k(s)}k(\bar s)
\]
où $\bar s$ est le point fermé de $\bar S$. Notons que la projection
\[
\tag*{(4.1.6)}
V_{\bar s}\longrightarrow V_s
\]
induit une \uline{équivalence} des topos étales \inserted{par} SGA4 \uline{VIII}. (car $k(s)$ est sép. clos et $k(\bar s)$ en est une extension radicielle), et on \deleted{\unsure{fera}} \inserted{pourra} donc identifier $V_s$ et $V_{\bar s}$.
\scanpage{145}
\mspage{V.17}
% archive sheet number: 144

Finalement, on trouve un diagramme commutatif \deleted{à carrés cartésiens}
\diagram{\begin{tikzcd}[ampersand replacement=\&,column sep=large,row sep=large]
\text{(4.2)} \& V_s\arrow[r,hook,"i"] \& V \& V_{\unsure{\eta|y}}\arrow[l,hook',"j"'] \\
\& V_{\bar s}\arrow[u,"h_s"]\arrow[r,hook,"\bar i"'] \& \bar V\arrow[u,"h"] \& V_{\bar\eta}\arrow[u,"h_\eta"']\arrow[l,hook',"\bar j"]
\end{tikzcd}},
dans lequel \deleted{les flèches verticales} (à isomorphisme de topos étales près tout au moins, en ce qui concerne $V_{\bar s}$) $V_{\bar s}$ et $V_{\bar\eta}$ sont resp. les images inverses des sous-schémas \inserted{fermé resp. ouvert} $V_s$ et $V_\eta$ de $V$, \deleted{et} et où $h_s$ \deleted{est} \inserted{définit} une équivalence de topos étales.

3.5. Les \deleted{\unsure{rôles des}} cinq schémas \inserted{essentiellement distincts} $V_s,V,V_\eta,\deleted{\unsure{V_{\bar s}}}V_{\bar\eta},\bar V$ jouent, du point de vue topologique, des rôles essentiels distincts :

\deleted{$\unsure{V_s|V_0}$} -- $V_s$ définit la structure topologique locale de la fibre $X_0$ en $x$

-- $V$ définit la structure topologique locale du schéma ambiant $X$ en $x$.

-- $V_\eta=V-V_s$ \deleted{\unsure{sert à relier}} \inserted{fournit, un lien entre,} les structures cohomologiques des schémas précédents $V_s$ et $V$

-- $V_{\bar\eta}$ fournit les \deleted{\unsure{cohomologies}} groupes des \deleted{\unsure{cohomologies}} cycles évanescents de $X$ en $x$. [\inserted{Il est} \deleted{\unsure{et est}} relié à $V_\eta$ $V_{\bar\eta}=V_\eta\times_\eta\bar\eta$, et sa cohomologie à celle de $V_\eta$ par la suite \deleted{exacte} spectrale habituelle].
\margin{par (4.2.1)\[
\left\{\begin{aligned}
(\underline\Phi^i_{\unsure{\nu|v|*}})_{\bar x}&\simeq H^{i-1}(V_{\bar\eta},\Lambda_\nu)\quad\text{si }i\geq2\\
0\longrightarrow(\underline\Phi^0_\nu)_{\bar x}&\longrightarrow\Lambda_\nu\longrightarrow H^0(V_{\unsure{\bar\eta|\eta}},\Lambda_\nu)\longrightarrow(\underline\Phi^1_\nu)_{\bar x}\longrightarrow0
\end{aligned}\right.
\]
suite exacte}

-- $\bar V$ permet de mettre en \unsure{relations|relation} les structures cohomologiques de $V_s\simeq V_{\bar s}$, et celle de $V_{\bar\eta}$.

On voit donc que les rôles de $V_\eta$ et $\bar V$
\scanpage{146}
\mspage{V.18}
% archive sheet number: 145

sont relativement \deleted{\unsure{fondamentaux}} \inserted{auxiliaires}, contrairement à \deleted{\unsure{celui|celle|ceux}} ceux de $V_s,V,V_{\unsure{\bar\eta|\bar y}}$. \deleted{\unsure{Il semble que}} On peut dire que la \deleted{\unsure{toute la structure|structure topologique}} \deleted{qui contiennent les informations} connaissance de la structure topologique de ces derniers contient \deleted{\unsure{structures topologiques des derniers|structures topologiques des fibres}} l'essentiel de \deleted{\unsure{ce que}} l'étude topologique de $f:X\longrightarrow S$ en le point $x$.

4.3. D'autre part, le groupe $I$ opère sur $\bar V$ \inserted{(3.4), via des opérations \unsure{sur}} par $V$-automorphismes, en induisant des opérations \uline{triviales} sur $V_{\bar s}$, et sur $V_{\bar\eta}$ les \inserted{$V_\eta$-}opérations évidentes sur l'écriture $V_{\bar\eta}\simeq V_\eta\times_\eta\bar\eta$.

4.4. Explicitons maintenant, par des triangles exacts dans des catégories dérivées convenables, les principales relations \deleted{\unsure{\&}} entre les cohomologies des cinq schémas envisagés, en choisissant des coefficients $\Lambda_{\unsure{\nu|v}}=\mathbb Z/\ell^{\nu+1}\mathbb Z$ comme d'habitude. Nous considérons $\mathbb R\unsure{\Gamma|\mathscr T}_{\deleted{\unsure{I|*}}}(\bar V,\Lambda_\nu)$ comme élément de $D(\Lambda_\nu[I])$, catégorie dérivée de la catégorie des $\Lambda_\nu$-modules à action admissible de $I$. Pour rappeler ce fait, nous écrirons plutôt $\mathbb R(\Gamma,I)(\bar V,\Lambda_\nu)$ ; l'application du foncteur oubli à cet objet redonne $\mathbb R\Gamma(\bar V,\Lambda_\nu)$. On a alors un triangle exact

\deleted{\unsure{(4.4.1)}} (A)
\diagram{\begin{tikzcd}[ampersand replacement=\&,column sep=large,row sep=large]
\& \mathbb R(\Gamma,I)(V_{\bar\eta},\Lambda_\nu)\arrow[dl,"{\unsure{\partial|\delta|)}}"']\&\\
\mathbb R(\Gamma,I)(V_{\bar s},\underline\Phi^\bullet_{V_{\bar s}})\arrow[rr] \& \& \mathbb R(\Gamma,I)(\bar V,\Lambda_\nu)\arrow[ul,"\bar j^*"]
\end{tikzcd}},
où $\underline\Phi^\bullet_{V_{\bar s}}=\mathbb R.\Gamma_{V_{\bar s}}(\Lambda_{\nu\bar V})$, peut aussi s'interpréter comme
\scanpage{147}
\mspage{V.19}
% archive sheet number: 146

la restriction à $V_s$ du complexe $\underline\Phi^\bullet_{\unsure{\nu|v}}$ \inserted{à opérations $I$} sur $X_0$ déjà \inserted{essentiellement} introduit dans II 1 (du moins par ses faisceaux de cohomologie, qui y \deleted{\unsure{étaient}} sont notés $\underline\Phi_\nu^{q\unsure{?|2|\varepsilon}}$). Par application du foncteur oubli, ce triangle devient

(A')
\diagram{\begin{tikzcd}[ampersand replacement=\&,column sep=large,row sep=large]
\& \mathbb R\Gamma(V_{\unsure{\bar\eta|\bar y}},\Lambda_\nu)\arrow[dl] \& \\
\mathbb R\Gamma(V_s,\underline\Phi^\bullet_\nu)\arrow[rr] \& \& \mathbb R\Gamma(\bar V,\Lambda_\nu)\arrow[ul,"\bar j^*"]
\end{tikzcd}}.

\deleted{Appliquant à (A) le foncteur $\mathbb R\Gamma I$, on trouve un triangle exact}

\deleted{Lorsque} Supposons que
\[
\tag*{(4.4.1)}
(\Phi^i_\nu)_{\bar x}=0\qquad\text{pour }i<m,
\]
alors \deleted{la suite exacte associée à (A') donne}
\[
\tag*{(4.4.2)}
H^i(V_s,\underline\Phi^\bullet_\nu)=0\qquad\text{pour }i<m,
\]
et la suite exacte associée à (A') donne
\[
\tag*{(4.4.3)}
H^i(\bar V,\Lambda_\nu)\xrightarrow{\bar j^i}H^i(V_{\bar\eta},\Lambda_\nu)
\quad\begin{aligned}&\text{injectif pour }i\unsure{\leq|\geq} m-1\\&\text{bijectif pour }i<m-1.
\end{aligned}
\]
% margin note written sideways
\margin{ce qui peut aussi s'écrire sous la forme
\[
\tag*{(4.4.4)}
H^i_{\bar x}(\unsure{\bar X|\bar V},\Lambda_\nu)\longrightarrow(\underline\Phi^i_\nu)_{\bar x}\quad\text{bijectif si }\deleted{\unsure{\leq}}i<m\deleted{\unsure{4}},
\]
compte tenu des hom. canoniques
\[
\tag*{(4.4.5)}
H^i_{\bar x}(\bar X,\Lambda_\nu)\xleftarrow{\sim}\varinjlim_m H^i_{\unsure{x|\bar x}}(\unsure{V(m)|X(m)},\Lambda_\nu)\unsure{\}|*}\longrightarrow(\underline\Phi^i_\nu)_{\bar x}
\]
déduits des hom. $\underset{\text{(4.1.2)}}{H^i(\bar V,\Lambda_\nu)}-H^i(V_{\bar\eta},\Lambda_\nu)$ déduits de l'inclusion $V_{\bar\eta}\longrightarrow\bar V$, \deleted{\unsure{et des}} et des relations analogues définissant les $H^i_{\bar x}(\bar X,\Lambda_\nu)$ en fonction des $H^{\unsure{i|i-1}}(\bar V,\Lambda_\nu)$.}

En particulier, si (4.4.1) est vrai pour $m=+\infty$, par \unsure{précision|prévision|positivité} si $x$ est pt \deleted{\unsure{isolé}} de l'ens. des points de non lissité de $f:X\longrightarrow S$, on trouve que $\mathbb R(\Gamma,I)(V_s,\Phi^\bullet_\nu)=0$, i.e. que \deleted{$\unsure{\mathbb R\Gamma I}$}
\[
\deleted{\unsure{\mathbb Rb}}\mathbb R(\Gamma,I)(\bar V,\Lambda_\nu)\mathrel{\overset{\bar j^*}{\unsure{-|\longrightarrow}}}\mathbb R(\Gamma,I)(V_{\bar\eta},\Lambda_\nu)
\]
est un isom., i.e.
\[
H^i(\bar V,\Lambda_\nu)\xrightarrow[\sim]{\bar j^*}H^i(V_{\bar\eta},\Lambda_\nu),\qquad\text{pour }\forall i,
\]
\unsure{nous avons} (4.4.5) un isom pour tout $i$
\scanpage{148}
\mspage{V.20}
% archive sheet number: 147

Dans ce dernier cas, \deleted{$\unsure{\bar V\longrightarrow}$} l'inclusion $V_{\unsure{\bar\eta|\bar y}}\longrightarrow\bar V$ s'apparente à un isomorphisme du pt de vue cohomologique (coeff. des $\Lambda_{\unsure{\nu|v}}$), et on peut dire que les cinq schémas envisagés \deleted{\unsure{se sont}} \deleted{\unsure{\&}} se réduisent à quatre, $V_s,V,V_{\unsure{\hat\eta|\bar\eta|\eta}},V_{\bar\eta}$ (on est débarrassé du seul $\bar V$ de ces schémas qui n'est pas noethérien !).

4.5. L'inclusion $V_{\bar s}\lhook\joinrel\longrightarrow\bar V$ définit de $\unsure{=|\simeq|\bar m}$ un homomorphisme
\[
\tag*{(4.5.1)}
\unsure{R|\mathbb R}(\Gamma,I)(\bar V,\Lambda_\nu)\longrightarrow\mathbb R(\Gamma,I)(V_s,\Lambda_\nu)
\]
qui s'insère dans un triangle exact

\deleted{\unsure{(4.5.2)}} (B)
\diagram{\begin{tikzcd}[ampersand replacement=\&,column sep=large,row sep=large]
\& \mathbb R(\Gamma,I)(V_{\unsure{\bar s|s}},\Lambda_\nu)\simeq\mathbb R\Gamma(V_s,\Lambda_\nu)^{\mathrm{triv}}\arrow[dl,"{\unsure{\partial|\delta|)}}"'] \& \\
\mathbb R(\Gamma,I)(\bar V,(\Lambda_\nu)_{V_{\bar\eta},V})\arrow[rr] \& \& \mathbb R(\Gamma,I)(\bar V,\Lambda_\nu)\arrow[ul,"\bar i^*"]
\end{tikzcd}}.

L'application du foncteur oubli à ce diagramme donne

(B')
\diagram{\begin{tikzcd}[ampersand replacement=\&,column sep=large,row sep=large]
\& \mathbb R\Gamma(V_s,\Lambda_\nu)\arrow[dl,"\partial"'] \& \\
\mathbb R\Gamma(\bar V,(\Lambda_\nu)_{V_{\bar\eta},V})\arrow[rr] \& \& \mathbb R\Gamma(\bar V,\Lambda_\nu)\arrow[ul,"\bar i^*"]
\end{tikzcd}}.

\deleted{Sous la condition (4.5.2)}
